
\chapter*{Practical Lab Manual}
\addcontentsline{toc}{chapter}{Practical Lab Manual}

\section*{Introduction to the Lab Manual}

This lab manual is designed to help students:
\begin{itemize}
    \item Learn programming by doing
    \item Practice real-life problem solving
    \item Prepare for practical examinations
\end{itemize}

Each experiment includes:
\begin{itemize}
    \item Aim
    \item Concept
    \item Algorithm
    \item Program
    \item Output
    \item Viva Questions
\end{itemize}

\section*{General Lab Rules}

\begin{itemize}
    \item Follow teacher instructions carefully
    \item Write programs neatly
    \item Test programs before submission
    \item Do not fear errors
\end{itemize}

\newpage
\section*{Experiment 1: Printing Messages}

\textbf{Aim:}  
To display a simple message using Python.

\textbf{Concept Used:}  
Print statement.

\textbf{Algorithm:}
\begin{enumerate}
    \item Start
    \item Display message
    \item Stop
\end{enumerate}

\begin{center}
\begin{tikzpicture}[node distance=2cm, auto]
    \node (start) [ellipse, draw, fill=green!20] {Start};
    \node (print) [io, below of=start] {Print "Hello..."};
    \node (stop) [ellipse, draw, fill=red!20, below of=print] {Stop};

    \draw [arrow] (start) -- (print);
    \draw [arrow] (print) -- (stop);
\end{tikzpicture}
\end{center}

\textbf{Program:}
\begin{verbatim}
print("Hello, welcome to Python Programming!")
\end{verbatim}

\textbf{Output:}
\begin{verbatim}
Hello, welcome to Python Programming!
\end{verbatim}

\textbf{Viva Questions:}
\begin{itemize}
    \item What does print do?
    \item Is Python case-sensitive?
\end{itemize}

\newpage
\section*{Experiment 2: Using Variables}

\textbf{Aim:}  
To store and display student information.

\textbf{Concept Used:}  
Variables and data types.

\textbf{Algorithm:}
\begin{enumerate}
    \item Store name in variable
    \item Store roll number
    \item Display values
\end{enumerate}

\textbf{Program:}
\begin{verbatim}
name = "Rahim"
roll = 10
print("Name:", name)
print("Roll:", roll)
\end{verbatim}

\textbf{Output:}
\begin{verbatim}
Name: Rahim
Roll: 10
\end{verbatim}

\textbf{Viva Questions:}
\begin{itemize}
    \item What is a variable?
    \item Why do we use variables?
\end{itemize}

\newpage
\section*{Experiment 3: Input from User}

\textbf{Aim:}  
To take input from the user.

\textbf{Concept Used:}  
Input function.

\textbf{Program:}
\begin{verbatim}
name = input("Enter your name: ")
print("Welcome", name)
\end{verbatim}

\textbf{Output:}
\begin{verbatim}
Enter your name: Ayesha
Welcome Ayesha
\end{verbatim}

\textbf{Viva Questions:}
\begin{itemize}
    \item What does input() return?
\end{itemize}

\newpage
\section*{Experiment 4: Decision Making}

\textbf{Aim:}  
To check pass or fail.

\textbf{Concept Used:}  
if-else statement.

\textbf{Program:}
\begin{verbatim}
marks = int(input("Enter marks: "))
if marks >= 33:
    print("Pass")
else:
    print("Fail")
\end{verbatim}

\textbf{Viva Questions:}
\begin{itemize}
    \item What is conditional statement?
\end{itemize}

\newpage
\section*{Experiment 5: Using Loops}

\textbf{Aim:}  
To print numbers from 1 to 10.

\textbf{Concept Used:}  
for loop.

\textbf{Program:}
\begin{verbatim}
for i in range(1, 11):
    print(i)
\end{verbatim}

\textbf{Viva Questions:}
\begin{itemize}
    \item Why do we use loops?
\end{itemize}

\newpage
\section*{Experiment 6: Functions}

\textbf{Aim:}  
To create a simple calculator function.

\textbf{Concept Used:}  
Functions.

\textbf{Program:}
\begin{verbatim}
def add(a, b):
    return a + b

print(add(5, 3))
\end{verbatim}

\textbf{Viva Questions:}
\begin{itemize}
    \item What is a function?
\end{itemize}

\newpage
\section*{Experiment 7: Lists}

\textbf{Aim:}  
To store multiple student marks.

\textbf{Concept Used:}  
List.

\textbf{Program:}
\begin{verbatim}
marks = [70, 80, 90]
print(marks)
\end{verbatim}

\textbf{Viva Questions:}
\begin{itemize}
    \item What is a list?
\end{itemize}

\newpage
\section*{Experiment 8: File Handling}

\textbf{Aim:}  
To save student data in a file.

\textbf{Concept Used:}  
File handling.

\textbf{Program:}
\begin{verbatim}
file = open("students.txt", "w")
file.write("Rahim\n")
file.close()
\end{verbatim}

\textbf{Viva Questions:}
\begin{itemize}
    \item Why is file handling important?
\end{itemize}

\newpage
\section*{Experiment 9: Object-Oriented Programming}

\textbf{Aim:}  
To create a student class.

\textbf{Concept Used:}  
Class and object.

\textbf{Program:}
\begin{verbatim}
class Student:
    def __init__(self, name):
        self.name = name

s = Student("Karim")
print(s.name)
\end{verbatim}

\textbf{Viva Questions:}
\begin{itemize}
    \item What is a class?
\end{itemize}

\newpage
\section*{Experiment 10: GUI Programming}

\textbf{Aim:}  
To create a simple GUI window.

\textbf{Concept Used:}  
Tkinter.

\textbf{Program:}
\begin{verbatim}
import tkinter as tk

window = tk.Tk()
window.title("My App")
window.mainloop()
\end{verbatim}

\textbf{Viva Questions:}
\begin{itemize}
    \item What is GUI?
\end{itemize}

\newpage
\section*{Final Practical Project}

\textbf{Project Title:} Student Record System

\textbf{Description:}  
Create a program to:
\begin{itemize}
    \item Add student data
    \item Save data to file
    \item Display saved data
\end{itemize}

\textbf{Evaluation Criteria:}
\begin{itemize}
    \item Logic – 40\%
    \item Correct Output – 30\%
    \item Code Neatness – 20\%
    \item Viva – 10\%
\end{itemize}

\section*{Conclusion}

Practice is the key to learning programming.
Every error is a lesson.

Keep practicing!